% GENERATED FILE. Edit canonical lessons, then run npm run generate:books.
% canonical-source-hash: fnv1a64:86367fe0f583a64e
% canonical-lessons: SA-C409-aksah, SA-C409-nemih, SA-C409-arah, SA-C409-bhramayati, SA-C409-vartma, SA-R409-review-402-409

\chapter{Wheels and Tracks}
\label{ch:sa-wheels-and-tracks}
\hlchaptermodality{409}

\begin{chapteropening}
\textbf{By the end of this chapter:} \emph{I can name the parts of a wheel and the ruts it leaves.}

Complete the last lesson of chapter 409: I can name the parts of a wheel and the ruts it leaves.
\end{chapteropening}

\section[\texorpdfstring{\sk{अक्षः} (akṣaḥ)}{अक्षः (akṣaḥ)}]{\sk{अक्षः} (akṣaḥ) — an axle}

\label{lesson:SA-C409-aksah}



\begin{quote}
Before the new one: say the Sanskrit for an elephant goad, then the Sanskrit for a stable.
\end{quote}

\subsection*{You'll want to know: \sk{अक्षः}}
\textbf{\sk{अक्षः}} — \emph{akṣaḥ} — \textquotedblleft{}an axle\textquotedblright{}.

\begin{quote}
\sk{शकटस्य} \sk{अक्षः} \sk{दृढः} \sk{अस्ति।}
\end{quote}

\begin{quote}
\emph{śakaṭasya akṣaḥ dṛḍhaḥ asti.}
\end{quote}

\begin{quote}
The cart's axle is strong.
\end{quote}

A masculine noun.

\begin{cousinweb}
\textbf{\sk{अक्षः}} (\emph{akṣaḥ}) is the word MW gives for an axle; MW gives no derivation within Sanskrit, so its origin is uncertain; MW compares Latin \emph{axis} and Greek \emph{axōn} (MW).
\end{cousinweb}

\begin{grammarlens}[title={What you've built}]
One more word for getting around.
\end{grammarlens}

\subsection*{Your turn}
\begin{itemize}
\raggedright
  \item \emph{Say it:} \emph{akṣaḥ}
  \item \emph{Say it:} \emph{akṣaḥ}, once more
  \item \emph{Say it:} \emph{śakaṭasya akṣaḥ dṛḍhaḥ asti.} — The cart's axle is strong.
  \item \emph{Recall:} say the Sanskrit for an elephant goad, then the Sanskrit for a stable, then say \emph{akṣaḥ} again
\end{itemize}

\begin{tcolorbox}[breakable,colback=teal!4,colframe=teal!35!black,title={Before you move on}]
What does \sk{अक्षः} mean? (\textquotedblleft{}An axle\textquotedblright{}.) Say it once more.
\end{tcolorbox}
\section[\texorpdfstring{\sk{नेमिः} (nemiḥ)}{नेमिः (nemiḥ)}]{\sk{नेमिः} (nemiḥ) — a wheel's rim}

\label{lesson:SA-C409-nemih}



\begin{quote}
Before the new one: say the Sanskrit for a stable, then the Sanskrit for an axle.
\end{quote}

\subsection*{You'll want to know: \sk{नेमिः}}
\textbf{\sk{नेमिः}} — \emph{nemiḥ} — \textquotedblleft{}a wheel's rim\textquotedblright{}.

\begin{quote}
\sk{रथस्य} \sk{नेमिः} \sk{दृढा} \sk{अस्ति।}
\end{quote}

\begin{quote}
\emph{rathasya nemiḥ dṛḍhā asti.}
\end{quote}

\begin{quote}
The chariot wheel's rim is strong.
\end{quote}

A feminine noun.

\begin{cousinweb}
\textbf{\sk{नेमिः}} (\emph{nemiḥ}) comes from the root √nam, \textquotedblleft{}to bend\textquotedblright{}, which also gives \emph{namati}, \textquotedblleft{}bows\textquotedblright{}; MW gives \emph{nemi} as the felly or rim of a wheel (MW).
\end{cousinweb}

\begin{grammarlens}[title={What you've built}]
One more word for getting around.
\end{grammarlens}

\subsection*{Your turn}
\begin{itemize}
\raggedright
  \item \emph{Say it:} \emph{nemiḥ}
  \item \emph{Say it:} \emph{nemiḥ}, once more
  \item \emph{Say it:} \emph{rathasya nemiḥ dṛḍhā asti.} — The chariot wheel's rim is strong.
  \item \emph{Recall:} say the Sanskrit for a stable, then the Sanskrit for an axle, then say \emph{nemiḥ} again
\end{itemize}

\begin{tcolorbox}[breakable,colback=teal!4,colframe=teal!35!black,title={Before you move on}]
What does \sk{नेमिः} mean? (\textquotedblleft{}A wheel's rim\textquotedblright{}.) Say it once more.
\end{tcolorbox}
\section[\texorpdfstring{\sk{अरः} (araḥ)}{अरः (araḥ)}]{\sk{अरः} (araḥ) — a spoke}

\label{lesson:SA-C409-arah}



\begin{quote}
Before the new one: say the Sanskrit for an axle, then the Sanskrit for a wheel's rim.
\end{quote}

\subsection*{You'll want to know: \sk{अरः}}
\textbf{\sk{अरः}} — \emph{araḥ} — \textquotedblleft{}a spoke\textquotedblright{}.

\begin{quote}
\sk{चक्रस्य} \sk{दश} \sk{अराः} \sk{सन्ति।}
\end{quote}

\begin{quote}
\emph{cakrasya daśa arāḥ santi.}
\end{quote}

\begin{quote}
The wheel has ten spokes.
\end{quote}

A masculine noun.

\begin{cousinweb}
\textbf{\sk{अरः}} (\emph{araḥ}) comes from the root √ṛ, \textquotedblleft{}to move, to fit\textquotedblright{}; MW gives \emph{ara} as the spoke of a wheel, from the Rigveda on (MW).
\end{cousinweb}

\begin{grammarlens}[title={What you've built}]
One more word for getting around.
\end{grammarlens}

\subsection*{Your turn}
\begin{itemize}
\raggedright
  \item \emph{Say it:} \emph{araḥ}
  \item \emph{Say it:} \emph{araḥ}, once more
  \item \emph{Say it:} \emph{cakrasya daśa arāḥ santi.} — The wheel has ten spokes.
  \item \emph{Recall:} say the Sanskrit for an axle, then the Sanskrit for a wheel's rim, then say \emph{araḥ} again
\end{itemize}

\begin{tcolorbox}[breakable,colback=teal!4,colframe=teal!35!black,title={Before you move on}]
What does \sk{अरः} mean? (\textquotedblleft{}A spoke\textquotedblright{}.) Say it once more.
\end{tcolorbox}
\section[\texorpdfstring{\sk{भ्रमयति} (bhramayati)}{भ्रमयति (bhramayati)}]{\sk{भ्रमयति} (bhramayati) — sets spinning; turns round}

\label{lesson:SA-C409-bhramayati}



\begin{quote}
Before the new one: say the Sanskrit for a wheel's rim, then the Sanskrit for a spoke.
\end{quote}

\subsection*{You'll want to know: \sk{भ्रमयति}}
\textbf{\sk{भ्रमयति}} — \emph{bhramayati} — \textquotedblleft{}sets spinning; turns round\textquotedblright{}.

\begin{quote}
\sk{बालः} \sk{चक्रं} \sk{भ्रमयति।}
\end{quote}

\begin{quote}
\emph{bālaḥ cakraṃ bhramayati.}
\end{quote}

\begin{quote}
The boy sets the wheel spinning.
\end{quote}

Past (he …): \sk{सः} \sk{चक्रं} \sk{भ्रमितवान्}, \emph{saḥ cakraṃ bhramitavān}.

\begin{cousinweb}
\textbf{\sk{भ्रमयति}} (\emph{bhramayati}) is built on the root √bhram, \textquotedblleft{}to wander, to turn\textquotedblright{}, which gives \emph{bhramati}, \textquotedblleft{}wanders\textquotedblright{}; MW gives the causative participle \emph{bhramita} as made to go round, whirled round (MW).
\end{cousinweb}

\begin{grammarlens}[title={What you've built}]
One more word for getting around.
\end{grammarlens}

\subsection*{Your turn}
\begin{itemize}
\raggedright
  \item \emph{Say it:} \emph{bhramayati}
  \item \emph{Say it:} \emph{bhramayati}, once more
  \item \emph{Say it:} \emph{bālaḥ cakraṃ bhramayati.} — The boy sets the wheel spinning.
  \item \emph{Recall:} say the Sanskrit for a wheel's rim, then the Sanskrit for a spoke, then say \emph{bhramayati} again
\end{itemize}

\begin{tcolorbox}[breakable,colback=teal!4,colframe=teal!35!black,title={Before you move on}]
What does \sk{भ्रमयति} mean? (\textquotedblleft{}Sets spinning; turns round\textquotedblright{}.) Say it once more.
\end{tcolorbox}
\section[\texorpdfstring{\sk{वर्त्म} (vartma)}{वर्त्म (vartma)}]{\sk{वर्त्म} (vartma) — a wheel track; a rut}

\label{lesson:SA-C409-vartma}



\begin{quote}
Before the new one: say the Sanskrit for a spoke, then the Sanskrit for sets spinning.
\end{quote}

\subsection*{You'll want to know: \sk{वर्त्म}}
\textbf{\sk{वर्त्म}} — \emph{vartma} — \textquotedblleft{}a wheel track; a rut\textquotedblright{}.

\begin{quote}
\sk{पंके} \sk{शकटस्य} \sk{गभीरं} \sk{वर्त्म} \sk{अस्ति।}
\end{quote}

\begin{quote}
\emph{paṅke śakaṭasya gabhīraṃ vartma asti.}
\end{quote}

\begin{quote}
There is a deep cart track in the mud.
\end{quote}

A neuter noun.

\begin{cousinweb}
\textbf{\sk{वर्त्म}} (\emph{vartma}) comes from the root √vṛt, \textquotedblleft{}to turn, to roll\textquotedblright{}, which also gives \emph{vartate}, \textquotedblleft{}is present\textquotedblright{}; MW gives \emph{vartman} as the track or rut of a wheel, a path, a road (MW).
\end{cousinweb}

\begin{grammarlens}[title={What you've built}]
One more word for getting around.
\end{grammarlens}

\subsection*{Your turn}
\begin{itemize}
\raggedright
  \item \emph{Say it:} \emph{vartma}
  \item \emph{Say it:} \emph{vartma}, once more
  \item \emph{Say it:} \emph{paṅke śakaṭasya gabhīraṃ vartma asti.} — There is a deep cart track in the mud.
  \item \emph{Recall:} say the Sanskrit for a spoke, then the Sanskrit for sets spinning, then say \emph{vartma} again
\end{itemize}

\begin{tcolorbox}[breakable,colback=teal!4,colframe=teal!35!black,title={Before you move on}]
What does \sk{वर्त्म} mean? (\textquotedblleft{}A wheel track; a rut\textquotedblright{}.) Say it once more.
\end{tcolorbox}
\section[(dialogue)]{Looking back: from Rest Houses to Wheels and Tracks}

\label{lesson:SA-R409-review-402-409}



\begin{quote}
Say the Sanskrit for a pilgrims' rest house and a hut of leaves.
\end{quote}

\subsection*{Your turn}
Say these aloud:

\begin{itemize}
\raggedright
  \item a pilgrims' rest house, a hut of leaves, sits down to rest, lies down to sleep, water for washing the feet
  \item a route, long, hard to get through, a traveller on foot, roams
  \item a long lonely road, a wilderness, travel through the country, takes cover, deserted
  \item swiftness, in a hurry, dashes off, last, arrives
  \item gets on, gets off, a crush of people, a queue, free space
  \item a driver of horses, a rein, a whip, a goad, harnesses
  \item a draught animal, a porter, a mahout, an elephant goad, a stable
  \item an axle, a wheel's rim, a spoke, sets spinning, a wheel track
\end{itemize}

\begin{tcolorbox}[breakable,colback=teal!4,colframe=teal!35!black,title={Before you move on}]
A pilgrims' rest house? (\textbf{\sk{धर्मशाला}}, \emph{dharmaśālā}.) A hut of leaves? (\textbf{\sk{पर्णशाला}}, \emph{parṇaśālā}.) Sits down to rest? (\textbf{\sk{निषीदति}}, \emph{niṣīdati}.) Lies down to sleep? (\textbf{\sk{संविशति}}, \emph{saṃviśati}.) Water for washing the feet? (\textbf{\sk{पाद्यम्}}, \emph{pādyam}.) A route? (\textbf{\sk{पन्थाः}}, \emph{panthāḥ}.) Long? (\textbf{\sk{दीर्घः}}, \emph{dīrghaḥ}.) Hard to get through? (\textbf{\sk{दुर्गमः}}, \emph{durgamaḥ}.) A traveller on foot? (\textbf{\sk{पदातिः}}, \emph{padātiḥ}.) Roams? (\textbf{\sk{विचरति}}, \emph{vicarati}.) A long lonely road? (\textbf{\sk{प्रान्तरम्}}, \emph{prāntaram}.) A wilderness? (\textbf{\sk{कान्तारः}}, \emph{kāntāraḥ}.) Travel through the country? (\textbf{\sk{देशाटनम्}}, \emph{deśāṭanam}.) Takes cover? (\textbf{\sk{निलीयते}}, \emph{nilīyate}.) Deserted? (\textbf{\sk{शून्यः}}, \emph{śūnyaḥ}.) Swiftness? (\textbf{\sk{जवः}}, \emph{javaḥ}.) In a hurry? (\textbf{\sk{सत्वरम्}}, \emph{satvaram}.) Dashes off? (\textbf{\sk{प्रधावति}}, \emph{pradhāvati}.) Last? (\textbf{\sk{अन्तिमः}}, \emph{antimaḥ}.) Arrives? (\textbf{\sk{आयाति}}, \emph{āyāti}.) Gets on? (\textbf{\sk{अधिरोहति}}, \emph{adhirohati}.) Gets off? (\textbf{\sk{अवतरति}}, \emph{avatarati}.) A crush of people? (\textbf{\sk{सम्मर्दः}}, \emph{sammardaḥ}.) A queue? (\textbf{\sk{पंक्तिः}}, \emph{paṅktiḥ}.) Free space? (\textbf{\sk{अवकाशः}}, \emph{avakāśaḥ}.) A driver of horses? (\textbf{\sk{यन्ता}}, \emph{yantā}.) A rein? (\textbf{\sk{वल्गा}}, \emph{valgā}.) A whip? (\textbf{\sk{कशा}}, \emph{kaśā}.) A goad? (\textbf{\sk{प्रतोदः}}, \emph{pratodaḥ}.) Harnesses? (\textbf{\sk{युनक्ति}}, \emph{yunakti}.) A draught animal? (\textbf{\sk{वाहः}}, \emph{vāhaḥ}.) A porter? (\textbf{\sk{भारवाहः}}, \emph{bhāravāhaḥ}.) A mahout? (\textbf{\sk{आधोरणः}}, \emph{ādhoraṇaḥ}.) An elephant goad? (\textbf{\sk{अंकुशः}}, \emph{aṅkuśaḥ}.) A stable? (\textbf{\sk{अश्वशाला}}, \emph{aśvaśālā}.) An axle? (\textbf{\sk{अक्षः}}, \emph{akṣaḥ}.) A wheel's rim? (\textbf{\sk{नेमिः}}, \emph{nemiḥ}.) A spoke? (\textbf{\sk{अरः}}, \emph{araḥ}.) Sets spinning? (\textbf{\sk{भ्रमयति}}, \emph{bhramayati}.) A wheel track? (\textbf{\sk{वर्त्म}}, \emph{vartma}.)
\end{tcolorbox}
